\section{例外曲線の移送に対する面積障害と相対 membrane}
\label{exm:section}

第\ref{tc:section}節の central exceptional curves を、
非零 slice へどの意味で比較できるかを明確にする。
Punctured curve の compact subsets に対する既存比較は成立するが、
正面積の閉曲線全体を symplectic surface として運ぶことには、より強い制限がある。
以下ではその制限を証明し、代わりに非零 slice 内で実際に得られる
symplectic disks と、その境界に残る作用を明示する。
これらの disks を相対 membrane と呼ぶ。
Torus fibration の saturated cap や complete corner が得られた、という意味ではない。

\subsection{Exact な smoothing には閉じた正面積曲線を移せない}

\begin{proposition}[Whole exceptional surface の局所移送障害]
\label{exm:exact-obstruction}
$(M,\omega=d\lambda)$ を exact symplectic manifold、
$(\Sigma,\sigma)$ を closed oriented surface とし、
$e=\int_\Sigma\sigma>0$ とする。
任意の smooth map $f:\Sigma\to M$ に対して
\[
 \int_\Sigma f^*\omega=0
\]
である。従って $f^*\omega=\sigma$ となる map、特に symplectic immersion は存在しない。
$\Sigma$ を symplectic submanifold として含む近傍を $M$ へ
symplectically 埋め込むこともできない。

第\ref{tc:section}節の chosen graph model では、各 $\tau\ne0$ について、
central exceptional sphere $\ell_i$ の全体を含む symplectic neighborhood を
$(\widetilde V_\tau,\Omega_\tau)$ へ埋め込むことは不可能である。
これは $\tau\to0$ で一様な比較だけに対する制限ではない。
\end{proposition}

\begin{proof}
Stokes の公式から
$\int_\Sigma f^*\omega=\int_\Sigma d(f^*\lambda)=0$ となる。
近傍の embedding があれば、その $\Sigma$ への制限が反例となる。

Chosen graph model の非零 fiber は $\CC^3$ と同型であり、
命題\ref{tp:radial}は全域の exact symplectomorphism
\[
 \Phi_\tau:
 (\widetilde V_\tau,\Omega_\tau,\lambda_\tau)
 \longrightarrow(\CC^3,\omega_{\rm std},\lambda_{\rm std})
\]
を与える。
他方、命題\ref{tc:resolution}と式\eqref{tc:kahler}から
$\int_{\ell_i}\widetilde\Omega=e_i>0$ である。
従って最初の主張を適用できる。
\end{proof}

この障害は primitive の一致を仮定していない。
また、source の正面積を縮小しても、target が exact である限り障害は残る。
第\ref{exc:section}節の whole $E$ を含む surface bridge を、
その symplectic form のまま非零 local fiber へ移す方式も同様に排除される。
一方、$E\setminus D$ の compact subsets に対する比較、
中央曲線を切って境界を持つ領域へ置き換える操作、
または全体の特異 fibration を別に作る操作は排除しない。

\begin{corollary}[全曲線上の形式の近似にも残る下限]
\label{exm:form-gap}
$\nu$ を $\Sigma$ 上の正の area form、$V=\int_\Sigma\nu$ とし、
二形式 $g\nu$ の norm を $\|g\nu\|_{\nu,\infty}=\sup_\Sigma|g|$ と定める。
命題\ref{exm:exact-obstruction}の任意の $f$ について
\[
 \|f^*\omega-\sigma\|_{\nu,\infty}\ge\frac eV.
\]
従って任意の smooth maps $f_\tau:\Sigma\to\widetilde V_\tau$ を使っても、
$f_\tau^*\Omega_\tau$ を $\sigma$ に $C^0$ 収束させることはできない。
\end{corollary}

\begin{proof}
$f^*\omega-\sigma=g\nu$ と置けば $\int_\Sigma g\nu=-e$ であり、
$e\le V\sup|g|$ となる。
\end{proof}

ここで否定しているのは、閉曲線全体上での pullback forms の収束である。
Resolved total space 内で集合として中央曲線へ近づくことや、
puncture を除いた compact sets 上の $C^\infty$ 収束とは異なる。

\subsection{Actual hypersurface の面積の量子化}

\begin{proposition}[Small exceptional area の全体への移送障害]
\label{exm:quantization}
$(X_\tau,\omega_\tau)$ が
\[
 [\omega_\tau]=\mathfrak a\mathsf H,
 \qquad \mathfrak a>0,
 \qquad \mathsf H\in H^2(X_\tau;\ZZ)
\]
を満たすとする。
任意の smooth map $f:\Sigma\to X_\tau$ に対して
\[
 \int_\Sigma f^*\omega_\tau
 =\mathfrak a\langle\mathsf H,f_*[\Sigma]\rangle
 \in\mathfrak a\ZZ.
\]
従って $e\notin\mathfrak a\ZZ$ の曲線 $(\Sigma,\sigma)$ は、
面積を保って $X_\tau$ へ symplectically 移せない。
特に $0<e<\mathfrak a$ の場合に適用できる。

さらに系\ref{exm:form-gap}の記号で
\[
 \|f^*\omega_\tau-\sigma\|_{\nu,\infty}
 \ge\frac{\operatorname{dist}(e,\mathfrak a\ZZ)}V.
\]
$e\notin\mathfrak a\ZZ$ なら、maps の homotopy classes を変えても
全曲線上で $f_\tau^*\omega_\tau\to\sigma$ という $C^0$ 近似は不可能である。
\end{proposition}

\begin{proof}
$f_*[\Sigma]$ は integral homology class なので、
$\mathsf H$ との pairing は整数である。
差の積分の絶対値は $\operatorname{dist}(e,\mathfrak a\ZZ)$ 以上であり、
前の系と同じ積分評価で norm の下限を得る。
\end{proof}

本稿では $\mathsf H=c_1(\mathcal O_{X_\tau}(1))$ である。
Chosen polarization $\mathfrak aH+\eps A$ の $A$ は central divisor に支持されるため、
非零 slice の class は $\mathfrak a\mathsf H$ である。
Ordered triple $i<j<k$ の二つの exceptional areas は
\[
 e_1=(j-i)\eps,\qquad e_2=(k-j)\eps.
\]
例えば既存の smallness 条件に
$0<\eps<\mathfrak a/(d-1)$ を加えれば、これらは $(0,\mathfrak a)$ に入る。
六次では $5\eps<\mathfrak a$ が十分である。
この場合、whole exceptional sphere を area-preserving に移す方式は、
exact な local chart の外へ出るだけでも救われない。
これは $c_1$ の負性や rational curves の非存在を使う議論ではなく、
面積が $\mathfrak a$ 以上の symplectic spheres を排除してはいない。

\subsection{Puncture に残る作用と反対向きの充填}

\begin{proposition}[Punctured sphere の boundary period]
\label{exm:puncture-period}
$E\simeq\PP^1$、$\omega_E=e\omega_{\rm FS}$、
$\int_E\omega_{\rm FS}=1$ とする。
有限座標 $Y$、$U=Y^{-1}$、$R=|Y|^2$ を用いると、
それぞれの chart で滑らかな primitives は
\[
 \lambda_B=e\frac{R}{1+R}\,d\theta_Y,
 \qquad
 \lambda_C=e\frac{|U|^2}{1+|U|^2}\,d\theta_U
\]
と取れ、overlap で
\[
 \lambda_B=\lambda_C-e\,d\theta_U
\]
を満たす。
$D_R=\{|Y|\le\sqrt R\}\subset E$ の正向き境界について
\[
 \operatorname{Area}(D_R)
 =\int_{\partial D_R}\lambda_B
 =e\frac{R}{1+R}\longrightarrow e.
\]
一般に、正面積 $e$ の closed surface を puncture した compact exhaustion
$\Sigma_\epsilon$ を exact manifold $(M,d\lambda)$ へ symplectically 移せば、
transported boundary の $\lambda$-period は $e$ へ収束する。
\end{proposition}

\begin{proof}
$d\theta_Y=-d\theta_U$ と $R=|U|^{-2}$ を代入すれば primitive の差を得る。
残りは表示の直接積分と Stokes の公式である。
\end{proof}

$\partial\Sigma_\epsilon$ は deleted disk の positive meridian と逆向きである。
従って puncture-positive の向きで書いた limiting period は $-e$ となる。
この period は smooth exact gauge $\lambda\mapsto\lambda+df$ では変化しない。
同じ primitive の下で、この boundary を面積が零へ向かう disk に潰す方式は成立しない。

\begin{corollary}[反対向きの positive disk に必要な面積]
\label{exm:opposite-disk}
$D$ を境界 $\gamma$ と面積 $e>0$ を持つ symplectic disk とする。
同じ exact manifold 内に、境界が $-\gamma$ となる positive symplectic disk $D'$ はない。
また命題\ref{exm:quantization}の $X_\tau$ 内で $0<e<\mathfrak a$ とすると、
そのような $D'$ が存在するなら
\[
 e+\operatorname{Area}(D')=\mathfrak a n,
 \quad n\in\ZZ_{>0},
 \qquad \operatorname{Area}(D')\ge\mathfrak a-e.
\]
ここでは boundaries が逆向きの maps として一致すると仮定する。
$D\cup D'$ の embedding や seam の smoothness は仮定しなくてよい。
\end{corollary}

\begin{proof}
Exact manifold では
$\operatorname{Area}(D')=\int_{-\gamma}\lambda=-e$ となる。
全体の場合、$D+D'$ は integral two-cycle であり、その面積は
$\mathfrak a n$ である。両 disk の面積が正なので $n\ge1$ となる。
\end{proof}

従って small exceptional area を持つ membrane の境界を、
反対向きの別の small positive disk によって閉じることもできない。
複数の境界を接続する cobordism や、共有作用・fiber decomposition を変更する手術は、
この系の仮定とは別に検討する必要がある。

\subsection{非零 slice に存在する二つの明示 membrane}

\begin{proposition}[Exceptional areas を境界作用に持つ actual disks]
\label{exm:disks}
第\ref{tp:section}節の exact radial map $\Phi_\tau$ を用い、
$D_e=\{\zeta\in\CC:\rho_\zeta\le e\}$、$\rho_\zeta=\pi|\zeta|^2$ と置く。
次は $(\widetilde V_\tau,\Omega_\tau)$ 内の smooth embedded symplectic disks である。
\begin{align}
 d_{1,\tau}(\zeta)
  &=\Phi_\tau^{-1}(\zeta,0,0),
  &\zeta&\in D_{e_1},\label{exm:first-disk}\\
 d_{2,\tau}^{r,\vartheta}(\zeta)
  &=\Phi_\tau^{-1}\left(
    e^{2\pi i\vartheta}\sqrt{\frac{r+e_1+\rho_\zeta}{\pi}},
    e^{-2\pi i\vartheta}\zeta,0\right),
  &\zeta&\in D_{e_2},\label{exm:second-disk}
\end{align}
ここで $r\ge0$ と $\vartheta\in\RR/\ZZ$ は固定する。
それぞれの pullback primitive は $\rho_\zeta d\theta_\zeta$ であり、
面積と正向き境界 period は、それぞれ $e_1,e_2$ である。
第二の disk の full moment image は
\[
 (P_x,P_y,P_z)=(r+e_1+q,q,0),\qquad 0\le q\le e_2
\]
となる。
\end{proposition}

\begin{proof}
第一の式は標準 coordinate disk の exact pullback である。
第二の式では第一座標の angle が一定で、第三座標は零である。
従って $\lambda_{\rm std}$ の pullback は第二座標からの
$\rho_\zeta d\theta_\zeta$ のみとなる。
平方根の中は $e_1>0$ により正であり、中心まで smooth である。
第二座標から $\zeta$ を復元できるので embedding であり、
pullback symplectic form は $d\rho_\zeta\wedge d\theta_\zeta$ である。
これは中心も含めた標準 area form を表す。
Exactness と境界積分から面積の主張を得る。Moment の式は直接計算である。
\end{proof}

これらは abstract な disk のみではなく、chosen graph form 上の明示した領域である。
第一の disk では $\rho_x\le e_1$、$y=z=w=0$、
第二では $\rho_x\le r+e_1+e_2$、$\rho_y\le e_2$、$z=w=0$ となる。
従って $r,e_i,|\tau|$ を小さくすれば、proper graph resolution の
exceptional chain を含む所定近傍へ入る。
ただし任意に指定した global form との relative comparison を、
これらの座標評価だけで主張するものではない。

Central sphere との関係も明示できる。
第一の curve $\ell_1\setminus\{p_\infty\}$ では
\[
 \zeta=X\sqrt{\frac{e_1}{\pi(1+|X|^2)}}
\]
が $D_{e_1}$ の内部との symplectic identification である。
第二の curve について $r>0$ とし、central product chart で
$x=e^{2\pi i\vartheta}\sqrt{r/\pi}$ を固定する。
$\{x\}\times(E\setminus D)$ 上で
\[
 \zeta=Y\sqrt{\frac{e_2}{\pi(1+|Y|^2)}}
\]
と置くと、式\eqref{exm:second-disk}の $\Phi_\tau$-image は
第\ref{tp:section}節の $\Psi_0$ の制限そのものである。
正実 $\tau$-ray ではそこでの phase normalization とも一致する。
一般の $\tau$ では、同節の第二座標の定数 phase を合わせればよい。
従って punctured sphere の各 compact subset に対する既存比較を、
この disk が boundary action $e_2$ を伴って明示的に延長している。

一方、$r=0$ では moment image は central $\ell_2$ の線分と一致するが、
$\vartheta$ を変えると異なる disks を得る。
そのため、式\eqref{exm:second-disk}を central $\ell_2$ 自身の lift と呼んではならない。
各 $D_{e_i}$ の境界 circle を quotient すれば、面積 $e_i$ の sphere を得る。
これは symplectic cut による compactification
\cite[\S1.1]{lerman-cut-local} であり、元の smooth slice 内への
closed symplectic sphere の embedding ではない。

\subsection{第三 residue と membrane action の一致}

\begin{proposition}[Actual normal line 上の明示した logarithmic input]
\label{exm:log-input}
$S=\operatorname{Bl}_0\CC^2$、$D=\{U=0\}$ とし、
第\ref{ct:degree-section}項の actual normal lines
\[
 N_x\simeq\mathcal O,
 \qquad N_T\simeq\operatorname{pr}^*\mathcal O_{\PP^1}(-1)
\]
を考える。
$N_x$ には flat connection、$N_T$ には curvature
$\kappa_T=\operatorname{pr}^*\omega_{\rm FS}$ の connection を選べる。
任意の $w_T>0$ に対して $S\setminus D$ 上に $\beta$ が存在し、
\[
 d\beta=\omega_e+w_T\kappa_T,
 \qquad
 \beta+(e+w_T)d\theta_U
 \quad\text{は $D$ の近傍へ smooth に延びる。}
\]
従って必要収支 $w_D=e+w_T$ はこの chosen input で実現され、
\[
 w_D-w_T=e
\]
は命題\ref{exm:disks}の第二 membrane の境界作用と一致する。
\end{proposition}

\begin{proof}
二つの charts で
\[
 \eta_B=\frac{|Y|^2}{1+|Y|^2}d\theta_Y,
 \qquad
 \eta_C=\frac{|U|^2}{1+|U|^2}d\theta_U
\]
と置く。すると $\eta_B=\eta_C-d\theta_U$ であり、
$d\eta_B=d\eta_C=\operatorname{pr}^*\omega_{\rm FS}$ である。
Actual normal coordinates の transition は $T=Uy$ なので、
\[
 \alpha_T=d\theta_T+\eta_B=d\theta_y+\eta_C
\]
が一つの smooth connection を定める。
これは $[\kappa_T]=-c_1(N_T)$ の規約に従い、$\int_E\kappa_T=1$ を満たす。

第\ref{exc:section}節の $Y$-chart の primitive を $\lambda_e$、
$U$-chart の smooth primitive を $\lambda_e^C$ と書く。
Euclidean terms は共通であり、projective terms から
$\lambda_e=\lambda_e^C-e\,d\theta_U$ となる。
そこで
\[
 \beta=\lambda_e+w_T\eta_B
\]
と置くと、微分の式を満たし、かつ
\[
 \beta=(\lambda_e^C+w_T\eta_C)-(e+w_T)d\theta_U
\]
となる。括弧内は $D$ へ smooth に延びる。

Punctured $E$ の正向き境界に沿う $\beta$-period は $e+w_T$ へ収束する。
そのうち normal connection に由来する $w_T\eta_B$ の period は $w_T$ へ収束し、
残る $\lambda_e$ の period は $e$ へ収束する。
これが所要の action と residue の関係である。
\end{proof}

この命題では actual normal line の上で connection を選んでいる。
任意に先に指定した connection を literal に保持する主張ではない。
また $\beta$ は $D$ を越えて smooth ではなく、
$\eta_B$ も全 $E$ 上の smooth primitive ではない。
従って第\ref{ct:coherent-family}項の global primitive 仮定や
flat transport の仮定が、これによって自動的に満たされるわけではない。
Logarithmic input の整合と、Liouville flow に適合する全 fibration の整合は区別される。

\paragraph{Complete corner に向けた帰結.}
中央の正面積曲線を、その symplectic neighborhood ごと非零 slice へ運ぶ工程は使えない。
代わりに punctured region を運び、残る正の boundary actions $e_i$ と
第三 divisor の residue を対向領域で処理する必要がある。
式\eqref{exm:first-disk}--\eqref{exm:second-disk}はそのための actual disks を指定するが、
境界は一般に既存 fibration の whole torus fibers ではない。
Shared actions、$T^2$-families、normal first jets、全 partial faces との一致を
追加で構成しなければならない。
全 exceptional core や自由境界のない fibration の一般非存在は、本節からは従わない。
